\section{Selective-explanation results and stopping decision}
\textbf{Outcome: no demonstrated benefit and no established novelty.} The complete first comparison changed anomaly decisions without yielding usable selective explanations. Its one repair then failed the final-label contract and stopped. No question learning is justified by these results.

\begin{center}\small\input{selective_v1/figures/summary.tex}\end{center}
\textit{The table preserves the originally frozen metrics. Its cited condition's single alert-plus-unresolved response copied the example string; it is not a useful resolved/unknown separation. Original JSON metrics are not applicable. Reference-word recall is a lexical diagnostic, not explanation accuracy.}

Both modified conditions produced exactly the same 24 binary decisions. Each preserved five of twelve original alerts and agreed with original on seventeen of twenty-four inputs. On annotation-included inputs they preserved four of seven original alerts; on gap inputs, one of five. There were no new alerts where original was normal. This is more conservative behavior, not a measured reduction in false positives. The label changes cannot be called improvements without correct input-level labels.

All 18 schema-valid manual explanations had empty observation lists. Of the 22 schema-valid cited explanations, 17 had empty observation lists and four of the remaining five copied the template's example claim. Only one contained nonplaceholder observation text, which remains semantically unverified. The single unresolved item also copied the template example. The first schema parser correctly checked structure but was insufficient as a useful-content check; a separate post-run audit exposes this distinction without changing the frozen analysis.

\begin{figure}[ht]\centering\includegraphics[width=\linewidth]{selective_v1/figures/selective_responses.pdf}
\caption{Completed first selective run. Left: paired binary output counts. Right: structural compliance versus populated/nonplaceholder observation text. Nonplaceholder text is not verified truth or evidence support.}
\end{figure}

\textbf{Repair failure.} The amended second run stopped after six calls on the two Explosion013 views. The cited gap response generated a fenced JSON explanation but omitted the required final Output field. No decision was imputed from its other text. The six-call record includes one invalid binary response, no capped generations, one valid structured explanation per modified condition, and no complete-run record. The valid manual gap explanation also supplied citations despite the instruction to leave its frame list empty. Two examples per condition cannot support a performance comparison. No further retry was made.

\textbf{Validation and compute.} The complete selective run used 72 calls in 190.54 seconds, including 173.62 generation seconds. All 24 paired tensor checks passed; frozen dependencies remained unchanged. Its original 24 labels and raw texts exactly reproduced the previous pilot. Eight transferred complete-run artifacts and seven partial-run artifacts matched remote SHA256. The six-call repair records 23.119 generation seconds. Its missing completion record means allocated elapsed time is unavailable; file-mtime differences were inconsistent and are not used to estimate it. Including the first abstention pilot, actual consumption is \textbf{126 of 192 allowed calls}, with 313.132 seconds of recorded generation across all three runs. Only Ada5 was used and is released; no A16 or Adas6/7 allocation occurred. Both completed runs plus the repair's 1,800-second enforced timeout fit within the aggregate one-GPU-hour ceiling.

\textbf{Research conclusion.} We have not demonstrated independent control of useful explanation specificity while preserving detection. We have demonstrated a failure of these particular instruction/output designs. This does not prove that InternVL2 or selective explanation in general cannot work. Because citation generation is already established, a formatting fix alone would also not create method novelty. Existing UCA references cannot certify all claims against sampled pixels. Stop this prompt direction's GPU compute and do not optimize on the consulted test sources. A subsequent research proposal would need a differentiated mechanism and stronger existing evidence-linked evaluation, potentially an already annotated benchmark; new human annotation is not logically mandatory.

\subsection{All paired binary decisions from the complete run}
Codes are 0 normal and 1 anomaly. ``Ann.'' denotes annotation-included, not verified anomalous. The partial repair is stored separately and is not substituted into this table.
\small\input{selective_v1/figures/all_cases.tex}\normalsize
